\section{Reproducibility}

The repository contains the original dataset, the complete Jupyter Notebook,
the generated figures, this report, a \texttt{requirements.txt} file, and a
short README with execution instructions. A reviewer can create a virtual
environment, install the pinned dependencies with

\begin{verbatim}
pip install -r requirements.txt
\end{verbatim}

and run all cells in \texttt{data\_workflow.ipynb} from top to bottom.

The notebook uses relative file paths, reusable functions, explicit cleaning
rules, and intermediate validation checks. The unmodified data are stored in
\texttt{df}, while the cleaned data are stored separately in
\texttt{df\_clean}. This separation preserves the original input and makes the
transformations easier to inspect.

Git was used to record progress through multiple commits. Development work was
performed on an additional branch before being integrated with the main branch.
The commit history documents the progression from repository setup through
data ingestion, cleaning, exploratory analysis, visualization, interpretation,
and reporting.

These practices support computational reproducibility by preserving the code,
data, software dependencies, and development history needed to repeat the
analysis \parencite{danchev2022}.